\documentclass[11pt,a4paper]{article}

%% PACHETE MATEMATICE
\usepackage{amsmath,amssymb,amsfonts}
\usepackage{amsthm}
\usepackage{mathpazo}
\usepackage{eulervm}
\usepackage{enumerate}
\usepackage{color}
\usepackage{graphicx}
\usepackage[all]{xy}
\usepackage{xcolor}
\usepackage{framed}
\usepackage[unicode]{hyperref}
\setcounter{MaxMatrixCols}{20}
\usepackage{multicol}
% \usepackage[romanian]{babel}


%% ASPECT
\linespread{1.05}
\usepackage[T1]{fontenc}
\usepackage[utf8x]{inputenc}
\usepackage[margin=2cm]{geometry}
% \usepackage[color]{showkeys}
\usepackage{anyfontsize}
\usepackage{titlesec}
\titleformat*{\section}{\large\bfseries}

\usepackage{fancyhdr}
\pagestyle{fancy}
\rhead{\small \color{gray} Notițe de Adrian Manea}
\lhead{\small \color{gray} ETTI-AM2, 2017---2018, M. Joița \& A. Niță}


\usepackage[autostyle = false, style = german]{csquotes}
\usepackage[romanian]{babel}
\newcommand\qq{\enquote}

%% COMENZILE MELE
\renewcommand*{\proofname}{Dem.:}
\newcommand\bb{\mathbb}
\newcommand\dr{\mathrm}
\newcommand\kal{\mathcal}
\newcommand\fk{\mathfrak}
\newcommand\op{\oplus}
\newcommand\ot{\otimes}
\newcommand\ds{\displaystyle}
\newcommand\id{\indent}
\newcommand\nid{\noindent}
\newcommand\seq{\subseteq}
\newcommand\sneq{\subsetneq}
\newcommand\speq{\supseteq}
\newcommand\spneq{\supsetneq}
\newcommand\rar{\rightarrow}
\newcommand\Rar{\Rightarrow}
\newcommand\xrar{\xrightarrow}
\newcommand\lar{\leftarrow}
\newcommand\Lar{\Leftarrow}
\newcommand\xlar{\xleftarrow}
\newcommand\lrar{\leftrightarrow}
\newcommand\Lrar{\Leftrightarrow}
\newcommand\ideal{\trianglelefteq}
\newcommand\ai{\text{ a.\^{i}. }}
\newcommand\sk{(\textit{Schi\c{t}\u{a}}) }
\newcommand\rhup{\rightharpoonup}
\newcommand\rhdn{\rightharpoondown}
\newcommand\lhup{\leftharpoonup}
\newcommand\lhdn{\leftharpoondown}
\newcommand\vid{\emptyset}
\newcommand\wh{\widehat}
\newcommand\ol{\overline}
\newcommand\NN{\mathbb{N}}
\newcommand\RR{\mathbb{R}}
\newcommand\ZZ{\mathbb{Z}}
\newcommand\QQ{\mathbb{Q}}
\newcommand\CC{\mathbb{C}}

%% STILURILE MELE
\newtheoremstyle{dotless}{}{}{\itshape}{}{\bfseries}{:}{ }{}

\theoremstyle{dotless}
\newtheorem{theorem}{Teorem\u{a}}
\newtheorem{proposition}{Propozi\c{t}ie}
\newtheorem{lemma}{Lem\u{a}}[section]
\newtheorem{corollary}{Corolar}
\newtheorem{exercise}{Exerci\c{t}iu}

\newtheoremstyle{dotlessdr}{}{}{}{}{\bfseries}{:}{ }{}
\theoremstyle{dotlessdr}
\newtheorem{example}{Exemplu}
\newtheorem{definition}{Defini\c{t}ie}
\newtheorem{remark}{Observa\c{t}ie}




\begin{document}
% \definecolor{labelkey}{RGB}{222,0,0}

\begin{center}
  {\large\textbf{Seminar 8.1 \\
      Recapitulare parțial}}
\end{center}

\vspace{2cm}

\section*{Ecuații diferențiale de ordin superior}

1. Să se integreze următoarele ecuații diferențiale:
\begin{enumerate}[(a)]
\item $ y^{\prime\prime\prime} = \sin x + \cos x $;
\item $ y\cdot y^{\prime\prime} - 2{y^\prime}^2 = 0 $ (omogenă, $ y' = z $);
\item $ y^{\prime\prime} + {y^\prime}^2 = \frac{2}{e^y} $ (autonomă, $ y' = p = p(y) $);
\item $ yy^{\prime\prime} - {y^\prime}^2 = 0 $ (autonomă);
\item $ (x^2 + 1)y^{\prime\prime} - 2xy^\prime + 2y = 0 $ ($ x > 0 $), știind
  că are o soluție particulară de forma $ y_p(x) = ax + b $;
\item $ y^{(4)} - 3y^{(3)} + 3y^{(2)} - y^\prime = 0 $;
\item $ y^{(4)} - 2y^{(3)} + 2y^{(2)} - 2y^\prime + y = 0$;
\item $ y^{(2)} + 10y' + 25y = 4e^{-5x} $;
\item $ y^{(3)} - y^{(2)} + y' - y = \cos x $;
\item $ y^{(2)} - y = 4e^x $;
\item $ x^2 y^{(2)} + 3xy' + 4y = 5x, x > 0 $ (Euler);
\item $ x^2y^{(2)} - xy' + y = x + \ln x, x > 0 $ (Euler);
\item $ 4(x + 1)^2 y^{(2)} + y = 0, x > -1 $ (Euler).
\end{enumerate}

%%%%%%%%%%%%%%%%%%%%%%%%%%%%%%%%%%%%%%%%%%%%%%%%%%%%%%%%%%%%%%%%%%%%%%

\section*{Sisteme diferențiale}

1. Să se determine soluția generală $ X = \begin{pmatrix} x(t) \\ y(t) \end{pmatrix} $
a sistemelor diferențiale:
\begin{enumerate}[(a)]
\item
  $
  \begin{cases}
    x'(t) &= x + 2y \\
    y'(t) &= 2x + 4y
  \end{cases}
  $
\item
  $
  \begin{cases}
    x'(t) &= 2x + y \\
    y'(t) &= 3x + 4y
  \end{cases}
  $
\end{enumerate}

%%%%%%%%%%%%%%%%%%%%%%%%%%%%%%%%%%%%%%%%%%%%%%%%%%%%%%%%%%%%%%%%%%%%%%

\section*{Teoria stabilității}

1. Scriind ecuațiile diferențiale sub forma unor sisteme, să se studieze
stabilitatea soluțiilor:
\begin{enumerate}[(a)]
\item $ x^{\prime\prime} + 3x' + 5 = f(t) $, cu $ f $ continuă pe $ [0, \infty) $;
\item $ x^{\prime\prime} + 9x = 0 $;
\item $ x^{\prime\prime} + {x'}^2 + x' + x = 0 $;
\item $ x^{\prime\prime\prime} + 10{x^{\prime\prime}} ^2 + 4x' - 8x^3 = 0 $.
\end{enumerate}

\emph{Indicație:} Notînd $ x = x_1 $ și $ x' = x_2 $, putem rescrie ecuațiile
sub forma unor sisteme, pe care le interpretăm matriceal.

\vspace{1cm}

2. Să se studieze stabilitatea spre $ \infty $ a soluțiilor sistemelor:
\begin{enumerate}[(a)]
\item
  $
  \begin{cases}
    x' &= -x + z \\
    y' &= -2y - z \\
    z' &= y - z
  \end{cases}
  $
\item
  $
  \begin{cases}
    x' &= x + 2y - 2z + t^2 \\
    y' &= 2x + y + 2z + t \\
    z' &= -2x + 2y + z
  \end{cases}
  $
\end{enumerate}

\vspace{1cm}

3. Să se studieze stabilitatea soluției nule a sistemului:
\[
  \begin{cases}
    x' &= x + y^2 - \cos x + 1 + 4y \\
    y' &= \sin y + (x + 1)^2 - \cos x
  \end{cases}.
\]


%%%%%%%%%%%%%%%%%%%%%%%%%%%%%%%%%%%%%%%%%%%%%%%%%%%%%%%%%%%%%%%%%%%%%%

\section*{Linii și suprafețe de cîmp}

1. Determinați liniile de cîmp pentru cîmpurile vectoriale:
\begin{enumerate}[(a)]
\item $ \vec{v} = x\vec{i} + y\vec{j} + xyz \vec{k} $;
\item $ \vec{v} = (xz + y)\vec{i} + (x + yz) \vec{j} + (1 - z^2)\vec{k} $;
\item $ \vec{v} = (x^2 + y^2)\vec{i} + 2xy\vec{j} - z^2 \vec{k} $.
\end{enumerate}

\vspace{1cm}

2. Scriind sistemele diferențiale autonome asociate, să se determine
soluția ecuațiilor liniare cu derivate parțiale de ordinul întîi:
\begin{enumerate}[(a)]
\item $ (y - z) \frac{\partial u}{\partial x} + (z - x)\frac{\partial u}{\partial y} + %
  (x - y) \frac{\partial u}{\partial z} = 0 $;
\item $ z(x - y) \frac{\partial u}{\partial x} + z(x + y) \frac{\partial u}{\partial y} + %
  (x^2 + y^2) \frac{\partial u}{\partial z} = 0 $;
\item $ (4z - 5y) \frac{\partial u}{\partial x} + (5x - 3z) \frac{\partial u}{\partial y} + %
  (3y - 4x)\frac{\partial u}{\partial z} = 0 $.
\end{enumerate}

%%%%%%%%%%%%%%%%%%%%%%%%%%%%%%%%%%%%%%%%%%%%%%%%%%%%%%%%%%%%%%%%%%%%%%

\section*{Ecuații cu derivate parțiale}

1. Să se determine soluțiile ecuațiilor cvasiliniare de ordinul întîi:
\begin{enumerate}[(a)]
\item $ (z - y)^2 \frac{\partial z}{\partial x} + xz \frac{\partial z}{\partial y} = xy $;
\item $ x(y^3 - 2x^3) \frac{\partial z}{\partial x} + y(2y^3 - x^3) %
  \frac{\partial z}{\partial y} = 9z(x^3 - y^3) $;
\item $ 2xu \frac{\partial u}{\partial x} + 2yu\frac{\partial u}{\partial y} = %
  u^2 - x^2 - y^2 $.
\end{enumerate}

\vspace{1cm}

2. Să se determine suprafața de cîmp a cîmpului vectorial
\[
  \vec{v} = yz\vec{i} + xz \vec{j} + xy\vec{k},
\]
care trece prin curba de ecuație:
\[
  \begin{cases}
    x^2 + y^2 &= 4 \\
    y &= 0
  \end{cases}.
\]

\vspace{1cm}

3. Să se determine suprafața de cîmp a cîmpului vectorial:
\[
  \vec{v} = 2xz\vec{i} + 2yz\vec{j} + (z^2 - x^2 - y^2)\vec{k},
\]
care conține cercul de ecuații:
\[
  \begin{cases}
    z &= 0 \\
    x^2 + y^2 - 2x &= 0
  \end{cases}.
\]

%%%%%%%%%%%%%%%%%%%%%%%%%%%%%%%%%%%%%%%%%%%%%%%%%%%%%%%%%%%%%%%%%%%%%%

\section*{Ecuații cu derivate parțiale de ordinul al doilea}

1. Să se aducă următoarele ecuații la forma canonică, precizînd natura lor:
\begin{enumerate}[(a)]
\item $ u_{xx} + u_{xy} - 12 u_{yy} = 0 $;
  \footnote{Am folosit notația prescurtată, %
    $ u_{xx} = \frac{\partial^2 u}{\partial x^2} $ etc.}
\item $ u_{xx} - 6u_{xy} + 9 u_{yy} = 0 $;
\item $ u_{xx} + 4u_{xy} + 13u_{yy} = 0 $.
\end{enumerate}

\vspace{1cm}

2. Rezolvați ecuațiile următoare, cu condițiile inițiale date:
\begin{enumerate}[(a)]
\item $ \begin{cases}
    u_{xx} - 4x^2 u_{yy} - \frac{1}{x} u_x &= 0, x > 0 \\
    u(1, y) &= y^2 + 1 \\
    u_x(1, y) &= 4
  \end{cases}
  $
\item $ \begin{cases}
    xu_{xx} - u_{yy} + \frac{1}{2} u_x &= 0, x > 0 \\
    u(x, 0) &= x \\
    u_y(x, 0) &= 0
  \end{cases}
$
\end{enumerate}



\end{document}